\advsection{Элементы тензорной алгебры.}
\advsubsection{Немного о ковекторах и их конкретном тензорном произведении}

\advsubsection{Тензорное произведение векторных пространств.}
Пусть $S$~---~произвольное множество.
Формальной линейной комбинацией его элементов будем называть функцию $f: S \to \R$ такую, что $f(s) = 0$ для всех элементов $S$ кроме, быть может, конечного числа.
\begin{definition}
    Свободным вещественным векторным пространством над множеством $S$ будем называть множество всех формальных комбинаций элементов $S$ с естественной структурой векторного пространства (то есть поточечное сложение двух функций и умножение на скаляр из $\R$) и обозначать как $\mathcal{F}(S)$.
\end{definition}
Для каждого элемента $x \in S$ существует функция $\delta_x \in \mathcal{F}(S)$ которая принимает $1$ в точке $x$ и $0$ во всякой другой.
Далее мы не будем делать различия между $x$ и $\delta_x$ и думать об $S$ как о подмножестве $\mathcal{F}(S)$.
Таким образом каждый элемент $f \in \mathcal{F}(S)$ может быть записан единственным образом как $f = \sum_{i = 1}^m a_i x_i$ где $x_j \in S$ такие, что $f(x_j) \neq 0$ и $a_j = f(x_j)$.
Значит $S$ является базисом для $\mathcal{F}(S)$ и $\mathcal{F}(S)$~---~конечномерное титк $S$~---~конечно.
\begin{proposition}[Характеристическое свойство свободного векторного пространства.]
    Для всякого множества $S$ и векторного пространства $W$ всякое отображение $A: S \to W$ единственным образом продолжается на линейное отображение $\overline{A}: \mathcal{F}(S) \to W$.
\end{proposition}
\begin{proof}
    Определим $\overline{A}$ на $f \in \mathcal{F}(S)$ следующим образом: поскольку $f = \sum\limits_{i = 1}^m a_{i}x_i$, то определим
    \[
        \overline{A}f = \sum\limits_{i = 1}^m a_{i}A(x_i).
    \]
    Определённое таким образом отображение очевидным образом линейно и в силу того, что любое другое отображение, продолжающее $A$, будет совпадать с $\overline{A}$ на базисных векторах, то оно единственно.
\end{proof}
Пусть теперь $V_1, \ldots, V_k$~---~вещественные векторные пространства. Рассмотрим $\mathcal{F}(V_1 \times \ldots \times V_k)$~---~множество всех конечных формальных линейных комбинаций $k$-кортежей вида $(v_1, \ldots, v_k)$ с $v_i \in V_i$.
Рассмотрим подпространство $\mathcal{R}$ свободного пространство, порождённое всеми элементами вида
\[
    (v_1, \ldots, a v_i, \ldots, v_k) - a(v_1, \ldots, v_i, \ldots, v_k),
\]
И такого вида:
\[
    (v_1, \ldots, v_i + v_i', \ldots, v_k) - (v_1, \ldots, v_i, \ldots, v_k) - (v_1, \ldots, v_i', \ldots, v_k).
\]
\begin{definition}
    Определим тензорное произведение $V_1, \ldots, V_k$, обозначаемое как $V_1 \otimes \ldots \otimes V_k$ следующим образом:
    \[
        V_1 \otimes \ldots \otimes V_k = \mathcal{F}(V_1 \times \ldots \times V_k)/\mathcal{R}.
    \]
\end{definition}
\advsubsection{Ковариантные и контравариантные тензоры на векторных пространствах.}
\advsubsection{Алгебра симметрических тензоров.}
\advsubsection{Алгебра кососимметрических тензоров и алгебра Грассмана.}


